\section{Introduction}\label{sec:introduction}
Let $P$ be a finite set of $n$ points in the plane, with no three
collinear, where $n$ is even. A pair of points is a \emph{halving pair}
if its line leaves exactly $(n-2)/2$ other points in each open halfplane.
Write $h(P)$ for the number of unordered halving pairs. Because no three
points are collinear, distinct pairs determine distinct lines; thus this
convention also counts halving lines.

\begin{theorem}\label{thm:halving-power}
There exist absolute constants $\varepsilon>0$, $C<\infty$, and $n_0$
such that, for every even $n\ge n_0$ and every $n$-point set
$P\subset\mathbb R^2$ with no three collinear,
\[
 h(P)\le C n^{4/3-\varepsilon}.
\]
\end{theorem}

The improvement is in the exponent and is uniform over all configurations.
The proof is nonquantitative: it supplies neither a numerical value of
$\varepsilon$ nor effective values of $C$ and $n_0$.

\subsection{The halving-line problem and earlier bounds}
A \emph{$k$-set} of a planar point configuration is a $k$-point subset
strictly separable from the remaining points by a line. The $k$-set
problem asks for the largest possible number of these subsets at a
specified size and rank. Halving pairs concern its middle rank and
measure the complexity of balanced line separations.

The early work of Lov\'asz and of Erd\H{o}s, Lov\'asz, Simmons and
Straus established the $O(n^{3/2})$ scale for halving lines and developed
the rotating-line and dissection-graph viewpoint
\cite{Lovasz1971,ErdosLovaszSimmonsStraus1973}.
Pach, Steiger and Szemer\'edi improved the upper bound by an
iterated-logarithm factor \cite{PachSteigerSzemeredi1992}.
Dey then proved the bound $O(n(k+1)^{1/3})$ for $k$-sets with
$1\le k\le n/2$, giving $O(n^{4/3})$ for halving lines
\cite{Dey1998}. His use of convex chains and graph crossings provides
the finite starting point of the present argument.

The lower bounds developed on a different scale.
Erd\H{o}s, Lov\'asz, Simmons and Straus constructed configurations with
$\Omega(n\log n)$ halving lines
\cite{ErdosLovaszSimmonsStraus1973}. T{\'o}th obtained
$n\exp(\Omega(\sqrt{\log n}))$ halving pairs, and Nivasch subsequently
simplified the construction and improved its exponential constant
\cite{Toth2001,Nivasch2008}. These bounds remain far below a fixed
power $n^{1+\delta}$. More recently, Alonso, L{\'o}pez and Rodrigo
improved a lower-order term in the upper bound while retaining the
exponent $4/3$ \cite{AlonsoLopezRodrigo2024}.

A related route through measures was developed by Pinchasi
\cite{Pinchasi2009}. A line avoiding the point set and leaving half
its points on each side gives a balanced cut. Pinchasi studied the
changes between such cuts ordered by slope and related endpoint-change
bounds to planar measure concentration. Our activity parameter also
counts points that change status during an interval and later return;
Section~\ref{sec:finite} supplies the reduction from that parameter
to endpoint changes, uniformly over all ranks.

\subsection{Rank sweeps and the stronger conclusion}
For the sweep arguments, call a configuration \emph{generic} if no three
points are collinear, its first coordinates are distinct, its determined
slopes are distinct, and no proper intersection of two determined
segments lies on a third determined segment. Section~\ref{sec:finite}
justifies reducing the halving-pair theorem to this case.

Fix a generic configuration and an integer rank $0\le k\le n$. For a slope $s$ distinct from the determined slopes, mark the
$k$ points with smallest values of
$y-sx$. As $s$ increases, the marked set changes by exchanging one
point for another. We call such a change a \emph{switch}. For an
interval of this sweep, its \emph{activity} is the number of distinct
points whose membership changes. At the middle rank, halving pairs
are exactly the switches. Write $S_k(P)$ for the total number of switches
as the slope ranges from $-\infty$ to $+\infty$.

\begin{theorem}[Uniform bound over all ranks]\label{thm:all-ranks}
There are absolute constants $\varepsilon>0$ and $C<\infty$ such that
for every integer $N\ge1$, every generic $N$-point planar configuration
$P$, and every integer $0\le k\le N$,
\[
 S_k(P)\le C N^{4/3-\varepsilon}.
\]
\end{theorem}

The uniformity in $k$ is needed even for the middle-rank conclusion:
restricting a block of switches to its participating points changes the
rank. At ranks $0$ and $N$ the marked set is constant and $S_k(P)=0$.

The same uniformity gives a bound throughout the planar $k$-set range.
For an arrangement $\mathcal L$ of nonvertical straight lines, let
$\Lambda_k(\mathcal L)$ be the closure of the points on its lines with
exactly $k$ lines strictly below them. Write
$\operatorname{comp}(\Lambda_k(\mathcal L))$ for the total number of
vertices and edge pieces of this polygonal chain.

\begin{corollary}[Planar $k$-sets and straight-line levels]
\label{cor:k-sensitive}
There are absolute constants $\varepsilon_0>0$ and $C_0<\infty$ such
that, for every integer $n\ge2$ and every integer $1\le k\le n/2$,
the number of $k$-sets of any $n$-point set $Q\subset\mathbb R^2$ with
no three collinear is at most
\begin{equation}\label{eq:k-sensitive}
 C_0 n(k+1)^{1/3-\varepsilon_0}.
\end{equation}
The same bound holds for $\operatorname{comp}(\Lambda_k(\mathcal L))$
for any arrangement $\mathcal L$ of $n$ nonvertical straight lines
with no parallel pair and no three concurrent.
\end{corollary}

The all-ranks theorem first bounds each level in an arrangement of a given
size. A shallow cutting covers a low level by $O(n/(k+1))$ cells, each
crossed by only $O(k+1)$ lines; within each cell, the level has a fixed local rank.
Applying the uniform bound in every cell gives the dependence on $k$ above.
The full argument, including the passage from levels to $k$-sets, appears
in Section~\ref{sec:k-sensitive-proof} after the main theorem is proved.

This corollary concerns planar point sets and complete straight lines
only; it asserts no corresponding saving for pseudolines or line segments.

\subsection{Proof strategy}

The central estimate concerns disjoint intervals of sweep steps with
activity at least $b$. The usual crossing argument bounds their number
by $O((n/b)^{4/3})$. We prove a uniform little-$o$ improvement of this
packing bound as $n/b$ tends to infinity. Choosing one sufficiently
large fixed scale then gives a contracting recurrence and the power
saving in Theorem~\ref{thm:halving-power}.

Here is the geometric mechanism behind the packing improvement.
For an interval whose endpoint marked sets differ, let $L$ be the points
marked only at its start and $G$ those marked only at its end. If
$|L|=|G|=a$, joining every point of $L$ to every point of $G$ gives a
complete bipartite exchange graph with $a^2$ edges. Convex and concave chain partitions
control its crossings. If the packing bound were attained to within
a constant factor, we could pass to limiting measures of vertices,
edges, and crossings after assigning points compact rank coordinates.
Each such coordinate is a smoothed fraction of a family of convex
chain graphs lying below the point.
Two separated time windows provide a pair of rank coordinates that
parametrizes the plane. The crossing counts imply that a positive
portion of edge activity has an absolutely continuous projection in
such a coordinate pair.

Compact rank coordinates alone do not control physical distances:
a small region in the rank square could represent a highly distorted
part of the original plane. We use the order of the sweep directions
to recover Euclidean control on a positive part of the limiting measure.
The coordinates belonging to other time windows become coordinatewise
monotone functions. Their almost-everywhere derivatives force three
transverse families of sweep lines to have three transverse limiting
directions. A local rigidity lemma for these three families prevents
collapse after a suitable linear normalization. Crucially, its limits
are taken first at a fixed small rank scale and then far enough along
the sequence of configurations. We preserve that order when passing
to a Hausdorff limit of coordinate graphs.

The crossing bound also makes edge endpoint mass absolutely continuous
with respect to vertex mass. Together with the absolutely continuous
edge projection in the rank square, this lets us select a portion with
positive original edge activity whose vertex projection has bounded
density there. The geometric control transfers this
vertex measure to the normalized Euclidean plane, where a disk of
radius $r$ has mass $O(r^{32/17})$, and hence a thin strip of width $r$
has mass $O(r^{15/17})$. In a fixed bounded region, the active vertices
of a narrow angular window lie either in three such strips or among
vertices accounted for by marked tokens crossing the boundary. A token
follows a marked point through the successive exchanges. Each
token crosses that boundary only a bounded number of times. Summing
these activity bounds contradicts the positive limiting edge mass.
The exponents $32/17$ and $15/17$ are auxiliary spatial exponents;
the final saving $\varepsilon$ comes from the separate fixed-scale
recurrence.

Section~\ref{sec:finite} proves the finite reductions and states the
uniform packing estimate. Section~\ref{sec:rank-limits} constructs the
rank coordinates and limiting measures. Sections~\ref{sec:differentiability}
and~\ref{sec:zoom} establish the directional and metric control, and
Section~\ref{sec:spatial} constructs the spatial measure.
Section~\ref{sec:boundary} proves the boundary estimate and packing
contradiction. Section~\ref{sec:recurrence} completes the finite induction
and derives the $k$-set and level bounds.

\paragraph{Background.}
We use weak compactness and regularity of finite Borel measures on
compact metric spaces, the Radon--Nikodym and Lebesgue decomposition
theorems and finite-measure differentiation
\cite[Chapter~1, Section~3]{Simon2014GMT}, almost-everywhere differentiability
of scalar Lipschitz functions, and planar separation by a Jordan curve.
The conditional-measure and Lipschitz-differentiability statements are
specified at their applications. The crossing, monotone-coordinate,
critical-image, three-foliation, and boundary estimates needed for the
new packing argument are proved in the paper.
